\section*{Entire original circuit argument (Section 5.5, pp.9:15--16)}
\begin{sourceresult}{Theorem 38}
Let $f$ be a monotone bivariate Laurent polynomial such that $\supp(f)$ is convexly independent and $|\supp(f)|=k$. Then $f$ requires monotone h.p.-circuit with $k/4$ gates.
\end{sourceresult}
Theorem 38 implies Theorem 2 via Lemma 28.
\begin{proof}
The lower bound is proved using the following ``progress'' measure. Given $A\subseteq\mathbb R^2$ and $\epsilon\in\{0,1\}$, let $A^\epsilon:=A$ if $\epsilon=1$ and $A^\epsilon:=\emptyset$ if $\epsilon=0$. Given $v\in\mathbb R^2$, let $v+A:=\{v\}+A$. Let $\mathcal A$ be a finite set of finite subsets of $\mathbb R^2$. For functions $\epsilon:\mathcal A\to\{0,1\}$ and $v:\mathcal A\to\mathbb R^2$, let
\[\mathcal A_{\epsilon,v}=\bigcup_{A\in\mathcal A}(v(A)+A)^{\epsilon(A)}.\]
Let
\[\mu(\mathcal A)=\max_{\epsilon,v}\{|\mathcal A_{\epsilon,v}|:\mathcal A_{\epsilon,v}\text{ is convexly independent}\}.\]
\par\medskip\noindent\textbf{Claim.} Let $\mathcal A'=\mathcal A\cup\{B\}$ and $A_1,A_2\in\mathcal A$. Then
\begin{align}
\mu(\mathcal A')&\le\mu(\mathcal A)+|B|,\tag{3}\\
\mu(\mathcal A')&\le\mu(\mathcal A)+2,&&\text{if }B=u+A_1\text{ for some }u\in\mathbb R^2,\tag{4}\\
\mu(\mathcal A')&\le\mu(\mathcal A)+4,&&\text{if }B=A_1\cup A_2,\tag{5}\\
\mu(\mathcal A')&\le\mu(\mathcal A)+4,&&\text{if }B=A_1+A_2.\tag{6}
\end{align}
\begin{proof}[Proof of Claim]
Inequality (3) is straightforward.
To prove (4), suppose that $\epsilon,v$ are such that $\mathcal A'_{\epsilon,v}$ is convexly independent. Suppose $\epsilon(A_1)=\epsilon(B)=1$ and $v(A_1)+A_1\ne v(B)+B$; otherwise we have $|\mathcal A'_{\epsilon,v}|\le\mu(\mathcal A)$. Then $(v(A_1)+A_1)\cup(v(B)+B)=\{v(A_1),v(B)+u\}+A_1$ is convexly independent. Since $|\{v(A_1),v(B)+u\}|=2$, by Lemma 37, $A_1$ has size at most 2. This means $\mu(\mathcal A')\le\mu(\mathcal A)+2$ by (3).
For (5), observe that $\mu(\mathcal A')\le\mu(\mathcal A\cup\{u_1+A_1,u_2+A_2\})$ whenever $u_1,u_2\ne0$ are distinct and apply (4) twice.
Finally, we prove (6). If $B=A_1+A_2$ is not convexly dependent, it contributes nothing to $\mu$. Assume that $B$ is convexly independent and $|A_1|\ge|A_2|>0$. By Lemma 37, either $|A_1+A_2|\le4$ or $|A_2|=1$. In the former case, $\mu(\mathcal A')\le\mu(\mathcal A)+4$ by (3). In the latter, $A_2=\{u\}$ for some $u$ and $\mathcal A'=\mathcal A\cup\{u+A_1\}$ and we can apply (4).
\end{proof}
Let us call a h.p.-circuit transparent, if every gate in the circuit computes a polynomial with convexly independent support. Given a circuit $\Psi$ and a node $u$, let $\supp(u)$ be the support of the Laurent polynomial computed by $u$. Let $\mathcal A_\Psi$ be the set $\{\supp(u):u\in\Psi\}$.
Using the Claim, we can show that whenever a transparent and monotone $\Psi$ has $s$ gates then $\mu(\mathcal A_\Psi)\le4s$. The proof is by induction. The induction base $s=1$ trivially holds. It remains to verify the induction step. Let $u$ be an output gate of $\Psi$. If $u$ is also an input gate, apply (3).
If $u=u_1+u_2$ then $\supp(u)=\supp(u_1)\cup\supp(u_2)$ and (5) completes the proof.
If $u=u_1\times u_2$ then $\supp(u)=\supp(u_1)+\supp(u_2)$ and (6) completes the proof.
Finally, consider a monotone circuit $\Psi$ for $f$ of minimal size $s$. No gate in the circuit computes the zero polynomial (unless $f$ itself the zero polynomial). The circuit is transparent because a monotone computation does not cancel monomials unless multiplying by zero, and because $+,\times$ can not ``undo'' convex dependence. This means that $\mu(\mathcal A_\Psi)\le4s$. On the other hand, since $\supp(f)$ consists of $k$ convexly independent points, we have $\mu(\mathcal A_\Psi)\ge|\supp(f)|=k$.
\end{proof}
